\paragraph*{04.03.03.02. Alle für das Projekt relevanten Verhaltensregeln und Berufsvorschriften identifizieren und einhalten}
\kcilabel{para:04.03.03.02_VerhaltensregelnUndBerufsvorschriften}{04.03.03.02}
\label{para:04.03.03.02_VerhaltensregelnUndBerufsvorschriften}


%%% Task
\par Als \gls{wsl} \gls{ricefw} identifizierte ich relevante Verhaltensregeln und Berufsvorschriften (Arbeitszeit, Arbeitsschutz, Betriebsvereinbarungen, Datenschutz, SAP-Standards, Code of Conduct) und überführte sie in ein Governance-Modell. Ziel: Einhaltung sicherstellen, Gesundheit schützen, Lieferfähigkeit erhalten.

%%% Action
\par Strukturierte Identifikation mit \gls{hr}, Rechtsabteilungen (\gls{rec}, \gls{jur}, \gls{jdr}), langjährigen Mitarbeitenden und \gls{pmo}. Implementierte Arbeitszeit-Governance mit Kernarbeitszeit Montag-Donnerstag 09:00-17:30. Ergänzte \gls{dod} um Compliance-Prüfpunkte. Einführte abgestuftes Kommunikationsmanagement. Dokumentierte Risiken im Register.

%%% Result
\par Alle Vorschriften eingehalten; keine Verstöße im Verlauf der Produktivsetzungen. Verbesserte Zusammenarbeit und Arbeitsumfeld. Prüfung von produktionsnahen Datennutzung erhöhte Sensibilisierung für Compliance-Risiken.

%%% Reflect
\par Frühzeitige Einbindung von \gls{hr}, Legal, IT-Security und \gls{br} beibehalten. Einführung automatisierter Dashboards und Standardisierung von Templates zur Effizienzsteigerung.

